\documentclass[11pt]{article}
\usepackage[T1]{fontenc}
\usepackage[utf8]{inputenc}
\usepackage{lmodern,microtype,amsmath,amssymb,amsthm,mathtools}
\usepackage[margin=1in]{geometry}
\usepackage[colorlinks=true,linkcolor=blue,citecolor=blue,urlcolor=blue]{hyperref}
\setlength{\emergencystretch}{3em}
\newtheorem{theorem}{Theorem}[section]
\newtheorem{lemma}[theorem]{Lemma}
\newtheorem{proposition}[theorem]{Proposition}
\providecommand{\tightlist}{\setlength{\itemsep}{0pt}\setlength{\parskip}{0pt}}
\newcommand{\R}{\mathbb R}
\author{ChatGPT-6 Astra \and Claude Fable 5 and Opus 5 \and Kaarel H\"anni \and OpenAI Codex \and Dmitry Vaintrob\thanks{Authors are listed in alphabetical order.}}
\date{Research draft, 21 October 2026}
\title{A quartic-width learning bound for Bayesian networks stable under internal noise}
\hypersetup{pdftitle={A quartic-width learning bound for Bayesian networks stable under internal noise},pdfauthor={ChatGPT-6 Astra; Claude Fable 5 and Opus 5; Kaarel Hanni; OpenAI Codex; Dmitry Vaintrob}}
\begin{document}
\maketitle
\begin{abstract}
We give an explicit sample-complexity bound for Gaussian Bayesian tanh
networks restricted by stability of clipped predictions under simultaneous
internal disturbances. For a layered Boolean target circuit of width $C$,
the sufficient number of examples is $O(C^4\log^9(2+C))$ when the other
parameters are fixed. The dependence on depth is polynomial, and the
statistical accuracy factor is $1/\varepsilon$. The guarantee holds at every
learner width above the teacher width in the strong-prior regime, and at
every actual finite-width Gaussian scaling with exponent $\gamma\ge2$.
The proof adapts an earlier pruning-and-Bernstein compression argument to
the clipped-score prior and supplies explicit Gaussian posterior thresholds.
It does not establish quadratic dependence on teacher width, or this
improved rate at mean-field scaling. The proof has independent AI-agent
mathematical review but is not Lean verified.
\end{abstract}

\section{Model and purpose}
A polynomial-depth learning theorem is more useful when the polynomial has
a manageable degree. This paper extracts an explicit quartic dependence on
Boolean teacher width for a particular Gaussian neural-network prior,
without strengthening its existing internal-noise support condition.
The proof uses the deterministic compressor in the earlier depth
report~\cite{depth-report}; that compressor is not new here. The contributions
are its adaptation to clipped-score stability, explicit rate bookkeeping,
and the finite-scale posterior conclusion.

Let $d\ge1$, $L\ge2$, $H=L-1$, and $\mathcal X=\{-1,1\}^d$.
The network has $H$ tanh hidden layers, all of width $N$, and a linear
readout:
\[
 h^1=\tanh(W^1\bar x),\qquad h^s=\tanh(W^sh^{s-1}),\qquad
 z_w=v^\top h^H,\qquad \bar x=(x,1).
\]
The independent Gaussian entries have variances $t^2/d$ in $W^1$,
$t^2/N$ in internal matrices, and $t^2/N^2$ in $v$, for $0<t\le1$.
Write
\[
 \kappa_N(w)=\frac dN\|W^1\|_F^2+
 \sum_{s=2}^H\|W^s\|_F^2+N\|v\|_2^2,
 \qquad E_v=N\|v\|_2^2.
\]
The scale $t=N^{(1-\gamma)/2}$ corresponds to exponent $\gamma$.

Disturbances are added after every hidden activation and at the output:
\[
 h^1_\eta=\tanh(W^1\bar x)+\eta^1,\quad
 h^s_\eta=\tanh(W^sh^{s-1}_\eta)+\eta^s,\quad
 z_w(x;\eta)=v^\top h^H_\eta+\eta^{\rm out},
\]
with $\|\eta\|_*^2=\sum_s\|\eta^s\|_2^2/N+|\eta^{\rm out}|^2$.
Let $T(u)=\max\{-1/2,\min\{1/2,u\}\}$ and fix $c>0$.
The prior conditions the raw Gaussian $\mu_{N,t}$ on the event
\[
 \mathcal R_{N,c}=\left\{w:
 |T(z_w(x;\eta))-T(z_w(x))|\le1/16\quad
 \text{for every }x,\ \|\eta\|_*\le
 r_c(w):=\frac{c}{\sqrt{1+HE_v(w)}}\right\}.
\]
The noisy hidden values need not lie in $[-1,1]$.
For iid inputs $X_i\sim D$, training further conditions on
$f(X_i)z_w(X_i)>1/2$. Denote this posterior by $Q_{S,N,t,c}$.
The risk is $R(Q)=\mathbb E_{w\sim Q,X\sim D}
\mathbf1\{f(X)z_w(X)\le1/4\}$, which bounds ordinary sign error.
No initial clean-output margin is required at unobserved inputs.

The target $f$ is computed by a layered Boolean circuit of width $C$ and
$H$ stages, using NAND, copy, and negation gates. Constant wires and any
compilation overhead count toward width; its output is one final wire.
Idle stages can be padded with copies. This is a Boolean-circuit comparator,
not an arbitrary real-weight neural implementation.

\section{An explicit learning theorem}
The target construction in Appendix~\ref{sec:theory-target-cylinder}
gives an explicit $D_*=D_*(d,L,C,c)\ge1$ such that, for every $N\ge C$,
a subset of $\mathcal R_{N,c}$ has correct margin on the entire cube and
raw Gaussian mass at least $e^{-ND_*/t^2}$. In particular, all posteriors
used below are well-defined.

For $K\ge1$, define
\begin{align*}
 r&=\min\{1/6,c/\sqrt{1+HK}\},&
 m&=\left\lceil256Kr^{-2}[d+\log(400L)]\right\rceil,\\
 B&=\sqrt{4mK},&
 u&=\frac1{48(d+1+Lm)(1+m(B+1))^L},\\
 P&=m(d+1)+(L-2)m^2+m,&
 J&=P\log(3+2B/u).
\end{align*}

\begin{theorem}\label{thm:main}
For $0<\varepsilon,\delta<1$, it suffices to take
\begin{equation}\label{eq:sample}
 n\ge\frac2\varepsilon[J+\log(1/\delta)].
\end{equation}
The following choices of $K$ and scales are valid:
\begin{enumerate}
\item Set $K=8D_*$. At every $N\ge C$, allow all scales satisfying
\[
 t^2\le\min\left\{1,
 \frac{ND_*}{(P_N/2)\log2+\log(2/\varepsilon)}\right\},
 \quad P_N=N(d+1)+(L-2)N^2+N.
\]
\item Set $K=8D_*+4(d+L)+4\log(2/\varepsilon)$. Allow every $N\ge C$
and $0<t\le N^{-1/2}$, including every exponent $\gamma\ge2$.
\end{enumerate}
With probability at least $1-\delta$ over the sample, $R(Q_{S,N,t,c})\le
\varepsilon$ simultaneously over the indicated widths and scales.
\end{theorem}

For the original strong-prior parameter $\sigma$, one has $t=\sqrt N\sigma$.
The first clause supplies a sufficiently small positive threshold at every
fixed width. The second clause concerns actual finite-variance priors.
The sample event is independent of width and scale in both clauses.

The code budget satisfies
\[
 J=O\!\left[
 \left(dLKr^{-2}(d+\log L)+L^2K^2r^{-4}(d+\log L)^2\right)
 \log(2+d+L+K+r^{-1})\right].
\]
The appendix's formulas give, with a universal constant,
\[
 D_*\le O\!\left((d+LC)
 [\log^2(2C)+\log(2L)+\log(2+d+C)]+(1+c^2)C\right).
\]
Consequently the sample bound is polynomial in depth and
$O(C^4\log^9(2+C))$ when the remaining parameters are fixed.
The exact sample formula retains the $1/\varepsilon$ factor and the
additional logarithmic dependence entering the finite-scale energy cutoff.

\section{Uniform compression on an energy sublevel}
Fix $w$ with $\kappa_N(w)\le K$. Put $k_s=\|W^s\|_F^2$ for internal
layers and $k_L=N\|v\|_2^2$, so their sum is at most $K$.
In every internal row, and in the readout, delete entries with magnitude
larger than
\[
 \frac{8\sqrt K\,\|\text{row}\|_2}{r\sqrt N}.
\]
Zero rows remain zero. Denote the pruned weights by bars. The total
absolute mass removed from a nonzero row is at most its squared norm
divided by the threshold. Therefore, on any vector in $[-1,1]^N$,
the internal pruning error has norm at most $r\sqrt{Nk_s/K}/8$;
the readout error is at most $r\sqrt{k_L/K}/8$.

Independently choose $m$ uniform indices $S^s_j$ with replacement at every
hidden layer. Keep the selected first-layer rows and multiply every
subsequent sampled source weight by $N/m$. The virtual full-width pass is
\begin{align*}
 \widetilde h^1&=\tanh(W^1\bar x),\\
 \widetilde h^s_i&=\tanh\!\left(\frac Nm\sum_{j=1}^m
 \bar W^s_{i,S^{s-1}_j}\widetilde h^{s-1}_{S^{s-1}_j}\right),\\
 \widetilde z&=\frac Nm\sum_{j=1}^m
 \bar v_{S^H_j}\widetilde h^H_{S^H_j}.
\end{align*}
Its stored coordinates form an ordinary width-$m$ network, including
when indices repeat or $m>N$.

Conditional on $S^1,\ldots,S^{s-2}$, the virtual state
$\widetilde h^{s-1}$ is independent of the fresh list $S^{s-1}$.
For one uniform source index $I$, put
$Y=N\bar W^s_{:,I}\widetilde h^{s-1}_I$. The variance and norm bounds are
\[
 \mathbb E\|Y-\mathbb EY\|_2^2\le Nk_s,
 \qquad \|Y-\mathbb EY\|_2\le16\sqrt{NKk_s}/r.
\]
Hilbert-space Bernstein~\cite{tarres-yao}, with tolerance
$r\sqrt{Nk_s/K}/8$, bounds failure by
\[
 2\exp\!\left(-\frac{3mr^2}{640K}\right)
 \le2\exp\!\left(-\frac{mr^2}{256K}\right).
\]
The readout calculation is identical without factors $\sqrt N$.
Zero-energy stages have zero error. A union bound over the $2^d$ inputs
and at most $L$ stages makes total failure at most $1/200$ by the chosen $m$.

On success, pruning and sampling together have normalized stage error at
most $r\sqrt{k_s/K}/4$. Define disturbances in the original network by
\[
 \eta^1=0,\qquad
 \eta^s=\widetilde h^s-\tanh(W^s\widetilde h^{s-1}),\qquad
 \eta^{\rm out}=\widetilde z-v^\top\widetilde h^H.
\]
Since tanh is 1-Lipschitz, the original disturbed network produces the
virtual states exactly and, simultaneously at every input,
\[
 \|\eta\|_*^2\le\sum_s\frac{r^2k_s}{16K}\le r^2/16.
\]
No product of downstream operator norms is required in this comparison.

The expected energy of the sampled network is at most $K$:
\[
 \frac dm\mathbb E\|W^1_{\rm comp}\|_F^2
   =\frac dN\|W^1\|_F^2,\qquad
 \mathbb E\|W^s_{\rm comp}\|_F^2\le\|W^s\|_F^2,\qquad
 m\mathbb E\|v_{\rm comp}\|_2^2\le N\|v\|_2^2.
\]
These identities use independent index lists in adjacent layers; pruning
can only decrease the corresponding expectations. Markov's inequality
gives compressed energy at most $4K$ with probability at least $3/4$.
Its intersection with the all-input noise event is nonempty.

If $w\in\mathcal R_{N,c}$, then $r\le r_c(w)$. Choose a sampled network
in that intersection. Its clipped output differs from the original's by
at most $1/16$ everywhere, and every weight has magnitude at most
$B=\sqrt{4mK}$. Round its weights to a fixed grid with coordinate error
at most $u$. A forward Lipschitz recurrence, using rounded weights at
most $B+1$, bounds the output change by
\[
 u(d+1+Lm)(1+m(B+1))^L\le1/48.
\]
Clipping is 1-Lipschitz. Thus the rounded network has clipped discrepancy
at most $1/12$ everywhere. There are at most $e^J$ rounded networks.
This finite class is independent of $N,t$, and of the original network.

\section{Generalization and the posterior energy tail}
Use the ramp $\psi$ which is one at signed clipped scores at most $1/3$,
zero at scores at least $5/12$, and linear in between. If a robust network
of energy at most $K$ fits every training margin, its associated rounded
network $h$ has $\psi(f(X_i)T(h(X_i)))=0$ on all training points.
Every original test-margin failure has compressed ramp loss one.
For any fixed $h$ with population ramp loss greater than
$\varepsilon/2$, the probability of zero empirical ramp loss is at most
$e^{-n\varepsilon/2}$. Indeed, the probability of positive loss is at least
its mean. A union bound over the finite class and \eqref{eq:sample} show
that every robust fitting network of energy at most $K$ has risk at most
$\varepsilon/2$, on one event of sample probability at least $1-\delta$.
This event is simultaneous over widths and applies to all posterior
mixtures supported on that energy sublevel.

Under the raw Gaussian, $N\kappa_N/t^2$ has the chi-square distribution
with $P_N$ degrees of freedom. Its exponential moment and the target
cylinder's mass give, uniformly over training samples,
\begin{equation}\label{eq:posterior-tail}
 Q_{S,N,t,c}\{\kappa_N>K\}
 \le\exp\!\left[-\frac{NK}{4t^2}+\frac{P_N}2\log2+
                      \frac{ND_*}{t^2}\right].
\end{equation}
For $K=8D_*$, the first scale threshold in Theorem~\ref{thm:main}
makes this at most $\varepsilon/2$.
If $t^2\le1/N$, then $t^2P_N/N\le d+L$. Multiplying the exponent by
$t^2/N$ bounds it by
\[
 -K/4+D_*+(d+L)(\log2)/2\le-\log(2/\varepsilon)
\]
for the second choice of $K$. Since $N/t^2\ge1$, the tail is again at
most $\varepsilon/2$. The retained-risk bound plus this discarded mass
proves the theorem.

\section{Scope and provenance}
The quartic rate was already implicit in the energy-compression argument
of the 17 October depth report~\cite{depth-report}. The present argument
adapts that proof to the clipped-score support and makes its finite
Gaussian consequence explicit. The target construction below is reproduced
from the robust-prior paper~\cite{robust-paper}, including work to which
Claude contributed. The review records accompanying this paper identify
the independent AI-agent mathematical checks; this paper is not Lean
verified and has not received a full human mathematical review.

At mean-field scaling, the expected internal Frobenius energy grows with
width; for $1<\gamma<2$, it grows as $N^{2-\gamma}$. The full-energy
argument above therefore does not give this rate in that range. This is
a limitation of the proof, not a lower bound on learning. The larger
entropy-compression proof in the earlier robust-prior paper covers those
scales, with a much worse explicit polynomial. Obtaining quadratic
teacher-width dependence under the same support condition remains open.

\appendix
\section{A robust target cylinder at every width}
\label{sec:theory-target-cylinder}

We give the probability bound used in the statistical proof. It uses a
robust digital implementation, not a claim that every implementation is
stable. Let $N\ge C$. Partition the neurons into $C$ balanced clone groups
of sizes $k_a\in\{\lfloor N/C\rfloor,\lceil N/C\rceil\}$, so
$C/2\le N/k_a\le2C$. Let $P$ project onto vectors constant on each clone
group. Fill unused logical slots with available constants or copies of an
existing Boolean wire. The ideal wire
vector $b^s(x)$ then has entries in $\{-1,1\}$, belongs to the range of
$P$, and has Euclidean norm $\sqrt N$.

Use gain $A$ in the tanh simulation of every Boolean gate. For example,
NAND on $u,v\in\{-1,1\}$ is the sign of $1-u-v$; copies and negations
use $u$ and $-u$. Bias constants in internal layers are supplied by counted
constant wires. The ideal signed preactivation of every gate is at least
$A$. In the clone implementation, each coefficient from a logical wire is
divided equally among that wire's clones. Set
\begin{align*}
 e_*&=\frac1{16\sqrt{2C}},&
 \tau_0&=\frac{e_*}{4},&
 b_L&=\sqrt{16\log9+8\log(4L)},\\
 A&=\max\{32b_L,2\log(8/e_*)\},&
 a&=\max\{2,256c\},&
 \rho&=\min\{1/4,A/[4(d+1)],Ae_*/(64\sqrt C)\}.
\end{align*}
The teacher readout averages the clones of the output wire with total
coefficient $a$.

In Gaussian coordinates $\sqrt d\,W^1$, $\sqrt N\,W^s$, and $Nv$, all
entries have variance $t^2$. Restrict each first-layer coordinate, each
internal row's projection onto every normalized clone-group direction, and
each readout coordinate to an interval of radius $\rho t$ about its teacher
value. The internal matrices then decompose as
\[
 W^s=W_*^s+R^s+E^s,\qquad
 R^s=R^sP,\quad \|R^s\|_{\rm op}\le\rho t\sqrt C,
 \quad E^sP=0.
\]
The remaining Gaussian bulk $E^s$ is independent of the constrained
coordinates. A $1/4$-net of the unit sphere has at most $9^N$ points,
and the usual two-net Gaussian bound gives
\[
 \Pr\{\|E^s\|_{\rm op}>b_Lt\}
 \le2\exp(2N\log9-Nb_L^2/8)\le\frac1{2L}.
\]
Thus all bulk norms are at most $b_Lt$ with conditional probability at
least $1/2$.

\paragraph{Finite disturbance robustness.}
Allow a separate RMS disturbance of size $\tau_0$ at each hidden layer and
scalar output disturbance at most $\tau_0$. This includes the entire
total-budget ball $\|\eta\|_*\le\tau_0$.
Write $e_s=\|h^s_\eta-b^s\|_2/\sqrt N$.
If $e_{s-1}\le e_*$, every clone-group mean differs from its ideal bit by
at most $e_*\sqrt{N/k_a}\le1/16$. At most three teacher coefficients,
including a constant coefficient, enter a gate, so its intended signed
preactivation remains at least $A/2$.
For $\epsilon\in\{-1,1\}$ and real $u$, the elementary global secant bound
\[
 \epsilon q\ge A/2
 \quad\Longrightarrow\quad
 |\tanh(q+u)-\epsilon|\le2e^{-A/2}+8|u|/A
\]
holds: for $|u|<A/4$ use saturation, and otherwise the linear term is at
least two. Since $E^sb^{s-1}=0$ and
$\|h^{s-1}_\eta\|_2/\sqrt N\le1+\tau_0$, it follows that
\[
 e_s\le2e^{-A/2}+
          (8/A)\rho\sqrt C(1+\tau_0)
          +(8b_L/A)e_{s-1}+\tau_0\le e_*.
\]
Each of the four terms is at most $e_*/4$.
The first layer initializes the same invariant: its coordinatewise
preactivation perturbation is at most $\rho(d+1)\le A/4$, and its
postactivation disturbance contributes at most $\tau_0$ in RMS norm.
Consequently the noisy signed readout is at least
\[
 f(x)z_w(x;\eta)
 \ge\frac{15a}{16}-\rho(1+\tau_0)-\tau_0>1
 \qquad(x\in\mathcal X,\ \|\eta\|_*\le\tau_0).
\]

The readout center satisfies $\sqrt N\|v_*\|\ge a\sqrt{C/2}$, and its
permitted perturbation has $\sqrt N\|v-v_*\|\le\rho$.
Therefore $E_v\ge a^2C/8$, whence
\[
 r_c(w)\le\frac{2\sqrt2c}{a\sqrt C}
       \le\frac1{64\sqrt{2C}}=\tau_0.
\]
All clean and noisy clipped scores on the calibrated ball equal $f(x)/2$.
The entire cylinder, with the bulk-norm event imposed, therefore belongs
to $\mathcal R_{N,c}$ and satisfies the training likelihood for every sample.
This also proves why every fixed $c>0$ is allowed: a larger calibration
requires only a polynomially larger readout coefficient.

\paragraph{Gaussian mass.}
There are $Nq$ constrained independent scalar coordinates, where
$q=d+2+(L-2)C$. The sum of their squared centers is at most $NB_0$, with
\[
 B_0=9dA^2+18(L-2)CA^2+2a^2C.
\]
These loose constants cover the fan-in-two gate implementation above.
For $Z\sim\mathcal N(0,t^2)$, $\rho\le1/4$, and any center $m$,
\[
 \Pr\{|Z-m|\le\rho t\}
 \ge\frac{2\rho}{e\sqrt{2\pi}}\exp(-m^2/t^2).
\]
Multiplication over the coordinates and the independent bulk event gives
\begin{equation}\label{eq:standalone-cylinder-mass}
 \mu_{N,t}(\text{target cylinder})
 \ge\exp(-ND_*/t^2),\qquad
 D_*=B_0+q\log\frac{e\sqrt{2\pi}}{2\rho}+\log2.
\end{equation}
Here $D_*\ge1$ and is bounded by a universal polynomial in $d,L,C,c$.
The construction and estimate hold at every integer width $N\ge C$.

\begin{thebibliography}{9}
\bibitem{depth-report}
ChatGPT-6 Astra and Claude (Fable 5 and Opus 5).
\emph{Depth dependence, internal robustness, and Bayesian neural-network learning}.
Research report, 17 October 2026, Section 3, Theorem 3.1.
\href{https://github.com/kaarelh/AI-works/tree/main/neural-network-generalization/works-in-progress/depth-supporting-notes/round1-depth-report}{Source and build instructions}.
\bibitem{robust-paper}
ChatGPT-6 Astra, Claude Fable 5 and Opus 5, Kaarel H\"anni, and Dmitry Vaintrob.
\emph{A polynomial-depth learning bound for NNBayes on robust networks}.
Research manuscript, 18 October 2026.
\href{https://github.com/kaarelh/AI-works/tree/main/neural-network-generalization/main-results/robust-polynomial-depth}{Source and build instructions}.
\bibitem{tarres-yao}
Pierre Tarr\`es and Yuan Yao.
\emph{Online learning as stochastic approximation of regularization paths:
optimality and almost-sure convergence}.
Appendix A, Corollary A.2, equation (A-1).
\url{https://arxiv.org/abs/1103.5538}.
\end{thebibliography}
\end{document}
